% portada.tex: portada y página legal del libro.
% La armadura es el puente Pratt del ejemplo C (capítulo 12), con cada barra en
% el color de su fuerza axial: tracción, compresión o sin fuerza.

\begin{titlepage}
\begin{tikzpicture}[remember picture, overlay]
  % Franja superior con el título
  \fill[principal] (current page.north west) rectangle
    ([yshift=-0.52\paperheight]current page.north east);
  \fill[acento] ([yshift=-0.52\paperheight]current page.north west) rectangle
    ([yshift=-0.535\paperheight]current page.north east);
  \node[anchor=north west, text width=0.78\paperwidth, align=left,
      text=white, font=\sffamily]
    at ([xshift=0.11\paperwidth, yshift=-0.16\paperheight]current page.north west)
    {{\fontsize{34}{40}\selectfont\bfseries\spaceskip=0.28em
      Análisis matricial\\[2pt]
      de armaduras\\[2pt] con Python\par}
     \vspace{0.9cm}
     {\Large Aprende el método de rigidez\\[2pt]
      construyendo tu propia calculadora\par}};

  % El puente del ejemplo C
  % 9 unidades de luz con escala 1.3: centrado, 5.85 cm a cada lado del eje
  \begin{scope}[shift={([xshift=-5.85cm, yshift=-0.79\paperheight]
      current page.north)}, scale=1.3, line cap=round]
    \foreach \k in {0,...,6}
      \coordinate (p\the\numexpr\k+1\relax) at (1.5*\k, 0);
    \foreach \k in {0,...,4}
      \coordinate (p\the\numexpr\k+8\relax) at (1.5*\k+1.5, 1.5);
    \foreach \i/\j/\tono in {
        1/2/traccion, 2/3/traccion, 3/4/traccion, 4/5/traccion,
        5/6/traccion, 6/7/traccion,
        8/9/compresion, 9/10/compresion, 10/11/compresion, 11/12/compresion,
        1/8/compresion, 12/7/compresion,
        2/8/traccion, 3/9/compresion, 4/10/neutro, 5/11/compresion,
        6/12/traccion,
        8/3/traccion, 9/4/traccion, 11/4/traccion, 12/5/traccion}
      \draw[\tono, line width=4.5pt] (p\i) -- (p\j);
    \foreach \i in {1,...,12}
      \fill[white] (p\i) circle (3.4pt);
    \foreach \i in {1,...,12}
      \draw[principal, line width=1.1pt] (p\i) circle (3.4pt);
  \end{scope}

  % Autor
  \node[anchor=south west, font=\sffamily\Large, text=principal]
    at ([xshift=0.11\paperwidth, yshift=0.08\paperheight]current page.south west)
    {Eddy Eder Sucapuca Cruz};
\end{tikzpicture}
\end{titlepage}

% ---------------------------------------------------------------------------
% Página legal (reverso de la portada)
\thispagestyle{empty}
\null\vfill
{\small\raggedright
\textbf{Análisis matricial de armaduras con Python}\\
Aprende el método de rigidez construyendo tu propia calculadora\par
\medskip
Eddy Eder Sucapuca Cruz\par
Versión preliminar, \the\year.\par
\bigskip
© \the\year{} Eddy Eder Sucapuca Cruz.\par
\medskip
El texto y las figuras de este libro se distribuyen bajo la licencia Creative
Commons Atribución-NoComercial-CompartirIgual 4.0 Internacional (CC BY-NC-SA
4.0): \url{https://creativecommons.org/licenses/by-nc-sa/4.0/deed.es}. Se
permite copiarlos, distribuirlos y adaptarlos sin fines comerciales, citando
al autor y compartiendo las obras derivadas con la misma licencia.\par
\medskip
El código de la calculadora y de los notebooks se distribuye bajo la licencia
MIT: puede usarse, modificarse e incorporarse en otros programas, también
comerciales, conservando el aviso de autoría.\par
\medskip
El código, los notebooks y la versión web del libro están en
\url{https://\textorepositorio}.\par
\bigskip
\textbf{Cómo citar este libro:}\\
Sucapuca Cruz, E. E. (\the\year). \emph{Análisis matricial de armaduras con
Python: aprende el método de rigidez construyendo tu propia calculadora}.
\url{https://\textorepositorio}\par
\bigskip
Compuesto con \LaTeX{} en las tipografías Libertinus y Source Code Pro. Todos
los resultados numéricos y los gráficos del libro se generan ejecutando el
código que se muestra.\par}
\cleardoublepage
